\scanpage{15}
\mspage{33}
l'injectivité de $\varphi^n$. Soit maintenant $F\longrightarrow G$ un effacement de $(R^nf'_*)g'^*$ avec $G$ de torsion, et soit $C=G/F$. Alors on peut compléter le diagramme $(*)$ par un $0$ au bout de la deuxième ligne, ce qui implique la surjectivité de $\varphi^n$. La dernière assertion de (iii) est immédiate à partir de la première.

\uline{Lemme(8.3):} Soit
\diagram{\begin{tikzcd}[column sep=small]
\deleted{\unsure{X}}&X\arrow[dr,"f"']&&\bar X\arrow[ll,"\pi"']\arrow[dl,"\bar f"]\\
&&S&
\end{tikzcd}}
un diagramme commutatif de morphismes propres. Supposons\deleted{\unsure{e}} que $\pi$ soit surjectif. Notons par un \deleted{sigle le}\inserted{``prime'' l'effet d'un} changement de base $g:S'\longrightarrow S$, et soient $g':X'\longrightarrow X$ et $\bar g:\bar X'\longrightarrow\bar X$ les morphismes \deleted{déduits}\inserted{canoniques}.

(i) \deleted{Si (5.1)(i) est vrai pour $(g,\bar f,.)$ et pour $(g',\pi,.)$ il l'est pour $(g,f,.)$.} \inserted{Assertion vide.}

(ii) Supposons que (5.1)(i) soit vrai pour $(g,f,.)$, $(g,\bar f,.)$ et $(g',\pi,.)$. Si (5.1(ii)) est vrai pour $(g,\bar f,.)$ \deleted{\unsure{et pour} $\unsure{(g',\pi,.)}$} il l'est pour $(g,f,.)$.

(iii) Si (5.1)(iii) est vrai pour $(g,\bar f,.)$ et pour $(g',\pi,.)$, il l'est pour $(g,f,.)$.

\deleted{Démonstration. (i): Supposons d'abord que le faisceau $F$ sur $X$ est de la forme $F=\pi_*G$. On a}
\mspage{14}
